\chapter{Experimental Procedure} 
\label{ch:Durchfuehrung}

\section{Measurement Method}
\label{sec:Messverfahren}
The experiment is based on the rotating crystal method for spectrometry. In this approach, a crystal holder equipped with the respective detection crystal (initially LiF, later NaCl) is positioned in the X-ray beam path. A Geiger-Müller counter tube serves as the detector for intensity-measured X-ray radiation. The actual measurement process is performed automatically by a control program on the computer. First, a safety check of the device is carried out, ensuring that the lead-glass window is closed before the high voltage is activated.

For recording the bremsstrahlung spectrum, the crystal is rotated over a specific angular range while, for each angular position, the count rate of detected photons is measured over a predetermined time period. To increase accuracy in determining the characteristic lines ($K_\alpha$ and $K_\beta$), after a coarse preliminary scan, detailed measurements are performed only in the regions of expected maxima. Here, the step size is reduced to capture the shape of the peaks precisely.

Another part of the experiment investigates the relationship between the count rate and the accelerating voltage of the X-ray tube at a constant detection angle. Here, the voltage is varied stepwise while the count rate is registered for each setting. Finally, the entire measurement sequence is repeated with an NaCl crystal to determine the lattice constant using the already determined wavelengths, and from this, the Avogadro number. All measured values are recorded digitally and documented directly.

\section{Experimental Setup}
\label{sec:Versucshaufbau}
The experimental setup consists of three main components: the X-ray source, the goniometric detection system, and the evaluation electronics. The X-ray tube features a water-cooled copper holder and a molybdenum anode, arranged inside the vacuum glass bulb opposite a heated tungsten cathode. A high-voltage source generates the electric field for accelerating the electrons.

The detection system includes a precision goniometer, on which the crystal holder and the counter tube are mounted movably. The rotation of the crystal and the simultaneous movement of the counter occur synchronously according to the Bragg angle (1:2 relation between crystal and counter rotation), so that the detector always captures the reflected beam. The counter tube is connected to an amplifier and a pulse counter, which digitizes the signals and forwards them to a computer. For safety, the entire beam path is embedded in a lead housing, accessible through a swivel lead-glass window. An internal safety switch ensures that the high voltage is switched off immediately as soon as the window is opened.

\subsubsection*{Sketch of the Experimental Setup}
\label{sec:Skizze}
% The sketch shows schematically the arrangement of X-ray tube, collimator, the rotatable crystal holder, the Geiger-Müller counter tube, and the connections to the control unit. The angles of incidence and reflection paths according to Bragg's Law are marked.

\begin{figure}[h]
    \centering
    % Insert image here, e.g.: \includegraphics[width=0.8\textwidth]{aufbau_skizze}
    \caption{Schematic setup of the X-ray spectrometer following the rotating crystal method.}
    \label{fig:aufbau}
\end{figure}